% paper_assets/limitations.tex
% Auto-generated scientific limitations disclosure for academic manuscripts
% Generated strictly from real-data metrics in results/*.json

\section{Scientific Limitations and Methodological Disclosures}
\label{sec:limitations}

In accordance with strict computational remote sensing reproducibility standards, all analytical findings, downscaled fields, and disturbance metrics presented in this study must be evaluated within the context of the following fundamental methodological constraints:

\subsection{Proxy-Based Validation in Crop Anomaly Detection}
In the absence of dense in-situ agronomic ground truth, the crop monitoring module (\texttt{ml/crop}) relies on an unsupervised learning paradigm.
\begin{enumerate}
    \item \textbf{Monsoon Cloud Penetration}: Sentinel-1 C-band SAR microwave penetration (5.405 GHz) expands valid observational coverage from 0.0\% to 0.0\% (+342.9\% relative gain), preventing observation blackouts.
    \item \textbf{Precipitation Sensitivity}: Optical NDVI correlates more strongly with CHIRPS regional rainfall deficits ($r = 0.576$) than the fused autoencoder anomaly ($r = 0.041$, $\Delta r = -0.536$). The fused autoencoder isolates fine-scale dielectric and structural canopy anomalies rather than linear macro-drought response.
    \item \textbf{Resolution Mixture}: Smallholder parcel boundaries ($<0.8\text{ ha}$ in Andhra Pradesh) introduce sub-pixel mixture effects across Sentinel $10\text{ m}$ resolution cells.
\end{enumerate}

\subsection{Hansen Global Forest Change Reference Label Realities}
Supervision and benchmarking for forest canopy disturbance (\texttt{ml/forest}) are derived from Hansen Global Forest Change (GFC v1.11, Landsat $30\text{ m}$).
\begin{enumerate}
    \item \textbf{Primary Test Region}: The Multimodal Siamese U-Net achieves $F_1 = 0.4858$ ($\text{IoU} = 0.3208$), outperforming Spectral Differencing ($\Delta\text{NDVI}~F_1 = 0.1739$, $\Delta\text{NBR}~F_1 = 0.1708$) and Torres et al. Early-Concat U-Net ($F_1 = 0.0204$).
    \item \textbf{Cloud Robustness}: Under heavy cloud cover, Sentinel-1 SAR integration enables the Siamese model to retain $F_1 = 0.0238$ (vs. $0.0204$ for optical-only Torres U-Net), though trailing simple NBR thresholding ($F_1 = 0.1288$).
    \item \textbf{Out-of-Region Domain Shift}: Evaluating models on Hasdeo Arand PEKB coal block yields an anticipated domain shift drop ($F_1 = 0.0000$), driven by Hansen temporal attribution lags.
    \item \textbf{Footprint Discrepancy}: Resampling $30\text{ m}$ Landsat loss masks onto $10\text{ m}$ Sentinel grids introduces boundary uncertainty (1 Hansen pixel = 9 Sentinel pixels).
\end{enumerate}

\subsection{Thermal Downscaling on Held-Out Spatial Blocks}
\begin{enumerate}
    \item \textbf{Spatial Generalization}: Under strict spatial block partitioning (1.92 km $\times$ 1.92 km blocks, 30\% spatial hold-out), Random Forest downscaling achieves $R^2 = 0.0608$, $\text{RMSE} = 2.194^\circ\text{C}$, and $\text{MAE} = 1.693^\circ\text{C}$. The modest out-of-sample $R^2$ is an expected physical consequence of complex coastal-urban terrain: 2D optical indices cannot resolve 3D urban canopy geometry, variable material thermal inertia, marine sea-breeze micro-advection, or point-source industrial plumes. Bounded RMSE demonstrates consistent regional calibration.
    \item \textbf{Skin vs. Ambient Temperature}: Landsat radiometric skin temperature ($T_{\text{skin}}$) differs by $3^\circ\text{C}$ to $12^\circ\text{C}$ from ambient air temperature ($T_{\text{air}}$).
    \item \textbf{Urban Forestry Deficits}: Target parcel cooling potentials (e.g. empirical -5.8$^\circ\text{C}$ deficit) are computed from satellite data relative to regional forest canopy refugia. Recommended tree species are literature-based suggestions, not computed from models.
\end{enumerate}
